\begin{figure}[htbp]
\centering
\begin{tikzpicture}[font=\small,>=Stealth]
  \node (pathpoint) at (0,0) {$n_i$};
  \node (oldpoint) at (3,1) {$a$};
  \node (probe) at (3,-1) {$m=p_0$};
  \node (commonmultiple) at (6,0) {$ae$};
  \draw[->] (pathpoint) -- node[above left] {$c_i$} (oldpoint);
  \draw[->] (pathpoint) -- node[below left] {$p_i$} (probe);
  \draw[->] (oldpoint) -- node[above right] {$e$} (commonmultiple);
  \draw[->] (probe) -- node[below right] {$b$} (commonmultiple);
\end{tikzpicture}
\caption{The relations defining the probe give two factorizations of
the prescribed common multiple $ae$.  Each arrow means right
multiplication by its label, not a covering relation.  This is not a
Hasse diagram: path points may coincide with $a$, and $a$ may divide
$m$.  The exact-profile assertion is that a divisor of $a$ omitted
from the path's downward closure has join $ae$ with $m$.}
\label{fig:probe-divisibility}
\end{figure}
